% Appendix "nn" (writer: oodcomp). Proofs from research/tracks/context/notes-final.md §4.3--4.4 (Prop 4.5, Thm 4.6,
% Cor 4.7, Rem 4.8); computations c2, c10, c12 in research/tracks/context/checks/.
\section{Proofs for Section~\ref{sec:nn}}\label{app:nn}

\begin{proof}[Proof of \cref{prop:nn:uniform}]
Let $C$ be a circuit for $f$ of size $s=\csize(f)<S$. The circuits of size $S$ whose first $s$ gates equal $C$, whose gates $s+1,\dots,S-1$ are arbitrary, and whose last gate is $\mathrm{AND}(\mathrm{gate}_s,\mathrm{gate}_s)$ all compute $f$. Their fraction among all circuits of size $S$ is $\prod_{j=s+1}^{S-1}16(d+1+j)^2\big/\prod_{j=1}^{S}16(d+1+j)^2=1/(N_s\cdot16(d+1+S)^2)$. For the mistakes: a law of a random total function has order-independent cumulative log loss $-\log P_{{\rm unif},S}(f)$ on every ordering (\cref{prop:setting:order}), and the expected number of Gibbs mistakes is at most $\ln2$ times the log loss, as in \cref{thm:comp:circuitbayes}(b).
\end{proof}

\begin{proof}[Proof of \cref{thm:nn:walsh}]
(a) Write $\mathcal K(x,x')=h(x\odot x')$ with $h(z):=\Psi(\sum_jz_j/d)$ and $\odot$ the coordinatewise product, using $\langle x,x'\rangle=\sum_jx_jx'_j$. For fixed $x$, $x'=x\odot z$ with $z$ uniform, and $\chi_S(x\odot z)=\chi_S(x)\chi_S(z)$, so $(T_{\mathcal K}\chi_S)(x)=\Ex_z[h(z)\chi_S(z)]\,\chi_S(x)$. The eigenvalue $\Lambda_S:=\Ex_z[h(z)\chi_S(z)]$ depends only on $|S|$, because $h$ is invariant under permutations of the coordinates. For the NNGP and NTK recursions with i.i.d.\ Gaussian weights and biases, every layer's kernel depends on the inputs only through $\langle x,x\rangle$, $\langle x',x'\rangle$ and $\langle x,x'\rangle$, and the first two equal $d$ on the cube.
(b) The characters form an orthonormal basis of $L^2$ of the uniform measure. For each $S$, $\min_{g\in\mathcal V}\|\chi_S-g\|^2=1-\|P_{\mathcal V}\chi_S\|^2$ with $P_{\mathcal V}$ the orthogonal projection, and $\sum_{\text{all }T}\|P_{\mathcal V}\chi_T\|^2=\operatorname{tr}P_{\mathcal V}=\dim\mathcal V\le n$. So $\sum_{|S|=k}\|P_{\mathcal V}\chi_S\|^2\le n$.
\end{proof}

\begin{proof}[Proof of \cref{cor:nn:separation}]
$\chi_S$ is total and computable, so \cref{thm:ood:cumulative} gives $\sum_te_t\le(\ln2)(K(\chi_S)+c)$ on every sequence of distinct inputs; repeated inputs have $e_t=0$, so the bound holds for i.i.d.\ inputs too (\cref{prop:setting:order}). An enumerative code of $S$ given $d$ and $k$ gives the bound on $K(\chi_S)$. Let $y^*=\chi_S(x)\in\{\pm1\}$ and let $w\le e$ be the posterior mass of the wrong label at $x$; the posterior mean is $\hat g=y^*(1-e-w)$, so $|y^*-\hat g|=e+w\le2e$ and $(y^*-\hat g)^2\le4e^2\le4e$. For i.i.d.\ inputs the expected test error after $t-1$ samples equals the expected online error at step $t$, which gives the averaged bound. The kernel side is \cref{thm:nn:walsh}(b) with $\mathcal V=\mathrm{span}\{\mathcal K(x_i,\cdot)\}$. Since $\log\binom dk\le k\log d$, the SFI side needs $n=O(k\log d/\varepsilon)$, while $\binom dk\ge(d/k)^k$.
\end{proof}

A direct computation ($d=10$, $20$ targets per cell, ridge $10^{-3}$; code in the accompanying repository) finds, for two- and three-layer ReLU NNGP, NTK and RBF kernels, squared errors above $1-n/\binom dk$ in every cell; for $k=5$ they stay above $0.68$ for all $n\le640$. A Bayes mixture over parities with prior $2^{-\ell(S)}$ has error $0.000$ from $n=20$ ($k=5$) and $n=40$ ($k=3$), and its cumulative Gibbs mistakes (mean $3.7$ to $5.5$, maximum $6.04$) stay below $(\ln2)\ell(S)$ ($6.70$ to $9.81$).

\begin{proof}[Proof of \cref{rem:nn:oddparity}]
(a) With $\arccos(-\rho)=\pi-\arccos\rho$ one gets $\Psi_1(\rho)-\Psi_1(-\rho)=\rho$ and $\rho\Psi_0(\rho)-(-\rho)\Psi_0(-\rho)=\rho$. So $\Psi_1(\rho)-\frac12\rho$ and $\Psi_1(\rho)+\rho\Psi_0(\rho)-\rho$ are even. The function $z\mapsto\Psi(\rho(z))-a\rho(z)$, $\rho(z)=\sum_jz_j/d$, is then invariant under $z\mapsto-z$, so its Walsh coefficients at odd $|S|$ vanish, while $a\rho(z)$ has Walsh coefficients only at $|S|=1$. Hence $\Lambda_S=0$ for odd $|S|\ge3$. Then $\langle\mathcal K(x_i,\cdot),\chi_S\rangle=(T_{\mathcal K}\chi_S)(x_i)=0$, so $\mathcal V\perp\chi_S$ and $\|\chi_S-\hat g\|^2=1+\|\hat g\|^2$.
(b) $\langle z,\chi_S\rangle:=\Ex_xz(x)\chi_S(x)$ is a linear functional of a centred Gaussian process, hence centred Gaussian, with variance $\Ex_{x,x'}\mathcal K(x,x')\chi_S(x)\chi_S(x')=\langle\chi_S,T_{\mathcal K}\chi_S\rangle=\Lambda_S=0$. So $\langle z,\chi_S\rangle=0$ almost surely. If $\mathrm{sign}\,z(x)=\chi_S(x)$ for all $x$, then $z(x)\chi_S(x)=|z(x)|$ and $\langle z,\chi_S\rangle=\Ex|z|>0$. So this event has prior probability $0$, and conditioning on an event of positive probability keeps it at $0$.
(c) From (a), $\mathcal K(x,x')=\sum_T\Lambda_T\chi_T(x)\chi_T(x')$, so $\sum_{x}\mathcal K(x_i,x)\chi_S(x)=2^d\Lambda_S\chi_S(x_i)=0$, and $\sum_{x}g(x)\chi_S(x)=0$ for every $g\in\mathcal V$; split the sum into training and unseen points. For kernel ridge regression $g=\sum_i\alpha_i\mathcal K(x_i,\cdot)$ with $\alpha=(K+\beta I)^{-1}y$ and $K$ the training Gram matrix, so $\sum_ig(x_i)y_i=y^\top K(K+\beta I)^{-1}y\ge0$: $K$ and $(K+\beta I)^{-1}$ are commuting positive semidefinite matrices. If some $g\in\mathcal V$ interpolates the training labels, its unseen sum is exactly $-n$.
\end{proof}

The computations behind \cref{rem:nn:oddparity} ($d=8$ and $d=10$; code in the accompanying repository): the eigenvalues at odd $|S|\ge3$ are below $3\cdot10^{-18}$ in absolute value for both bias-free two-layer kernels and nonzero with a bias or a third layer; the variance of $\langle z,\chi_S\rangle$ is at most $4.4\cdot10^{-18}$ for $|S|\in\{3,5\}$ against $1.5\cdot10^{-4}$ and $1.4\cdot10^{-3}$ at $|S|=4$; the identity of (c) holds to relative error $6.2\cdot10^{-12}$, the training sum is at least $89.7$ in every run, and the mean unseen product falls from $-0.10$ ($n=100$) to between $-0.81$ and $-0.92$ ($n=900$). The sign accuracy of kernel ridge regression on unseen points ($d=10$, ten targets per cell, ridge $10^{-3}$ times the mean diagonal) is $0.46$--$0.51$, $0.37$--$0.41$, $0.18$--$0.20$ and $0.05$--$0.07$ for $n=100,300,640,900$, while controls with a nonzero eigenvalue ($|S|=4$, or a three-layer NNGP with $|S|=3$) reach $0.99$--$1.00$ at $n=640$.
